\chapter{电动力学概念骨架：场、波与规范}
\section{Maxwell 方程}
\[
\begin{array}{c|c}
\text{微分形式}&\text{积分形式}\\ \hline
\nabla\cdot\vb E=\rho/\varepsilon_0&\oint\vb E\cdot\dd\vb a=Q/\varepsilon_0\\
\nabla\cdot\vb B=0&\oint\vb B\cdot\dd\vb a=0\\
\nabla\times\vb E=-\partial_t\vb B&\oint\vb E\cdot\dd\vb l=-\dd\Phi_B/\dd t\\
\nabla\times\vb B=\mu_0\vb J+\mu_0\varepsilon_0\partial_t\vb E&
\oint\vb B\cdot\dd\vb l=\mu_0I+\mu_0\varepsilon_0\dd\Phi_E/\dd t
\end{array}
\]
四式依次表达电荷是电场源、无磁单极、变化磁场产生涡旋电场、传导电流与变化电场产生磁场。位移电流项还保证
$\partial_t\rho+\nabla\cdot\vb J=0$。

\section{静场、Green 函数与镜像法}
静电场 $\vb E=-\nabla\phi$，代入 Gauss 定律：
\[
\nabla^2\phi=-\rho/\varepsilon_0,
\]
无源区即 Laplace 方程。自由空间 Green 函数满足
$\nabla^2G(\vb r,\vb r')=-\delta(\vb r-\vb r')$，
$G=1/(4\pi|\vb r-\vb r'|)$，从而
$\phi=\varepsilon_0^{-1}\int G\rho\,\dd^3r'$ 加边界项。接地无限导体平面上方的点电荷，可在镜像位置放置 $-q$；上半空间边界条件相同，唯一性定理保证所得场正确，但镜像电荷不是真实源。

静磁场因 $\nabla\cdot\vb B=0$ 可写 $\vb B=\nabla\times\vb A$。取 Coulomb 规范
$\nabla\cdot\vb A=0$，由 $\nabla\times(\nabla\times\vb A)=\mu_0\vb J$ 得
\[
\nabla^2\vb A=-\mu_0\vb J.
\]

\section{电磁波、能流与规范}
真空无源区对 Faraday 定律取旋度：
\[
\nabla\times\nabla\times\vb E=\nabla(\nabla\cdot\vb E)-\nabla^2\vb E
=-\mu_0\varepsilon_0\partial_t^2\vb E,
\]
因 $\nabla\cdot\vb E=0$，
\[
\boxed{\nabla^2\vb E-\frac1{c^2}\partial_t^2\vb E=0},\qquad
c=(\mu_0\varepsilon_0)^{-1/2}.
\]
平面波代入散度方程得到 $\vb k\cdot\vb E=\vb k\cdot\vb B=0$，且
$\vb B=(1/\omega)\vb k\times\vb E$，所以真空波横向。偏振描述 $\vb E$ 端点轨迹；能流为
$\vb S=\vb E\times\vb H$。

势表示
\[
\vb B=\nabla\times\vb A,\qquad
\vb E=-\nabla\phi-\partial_t\vb A
\]
在变换
$\vb A'=\vb A+\nabla\chi,\ \phi'=\phi-\partial_t\chi$ 下不变。规范自由度是描述冗余而非新物理；选 Coulomb 或 Lorenz 规范是方便计算。四维记号中
$J^\mu=(c\rho,\vb J)$，$A^\mu=(\phi/c,\vb A)$，
$F^{\mu\nu}=\partial^\mu A^\nu-\partial^\nu A^\mu$；Maxwell 方程压缩为
$\partial_\mu F^{\mu\nu}=\mu_0J^\nu$ 和 Bianchi 恒等式。

\pitfall{电磁波横向结论这里限定真空均匀无源区；波导中可有纵向分量。规范势可变不表示势毫无意义，Aharonov--Bohm 效应体现其量子相位作用。}
\section{核心公式与面试回答}
\interview{推导波动方程时务必依次说“取旋度—用矢量恒等式—使用无源散度方程—代入另一条旋度方程”。回答规范不变性时直接代入势变换，展示二阶混合偏导抵消。}
